\section{Shortest channels in arbitrary periodic configurations}
\label{sec:g-localization}
The relevant localization is a consequence of separated strict
convexity, not proximity to a particular model shape.

\begin{lemma}[A geometric criterion for the complete event]
\label{lem:g-channels}
At any reference configuration of Section~\ref{sec:geometric-results},
the minimum distance between distinct lifted obstacles is positive and
is attained by finitely many oriented pairs modulo common translation.
Every such pair has a unique closest segment, positive-definite distance
Hessian and positive clearance from the remaining obstacles. Its contact
points, length and curvatures continue smoothly under small smooth
deformations. There are fixed contact and direction neighborhoods and
$\varepsilon_0>0$ such that every physical sequence with $j$ complete
flights and total length at most $jg_*+\varepsilon_0$ alternates in one
of these neighborhoods. The constants are independent of $j$.
\end{lemma}
\begin{proof}
Choose bounded representatives of the finitely many obstacles. The
distance between $D_a$ and $D_b+\lambda$ tends to infinity with
$|\lambda|$, uniformly in $a,b$. Only finitely many lattice translates
can have distance below any fixed bound. Each of those finitely many
positive distances is attained; disjoint closures make their minimum
positive. Among pairs not attaining this minimum, the next distance is
separated from it by a positive amount. Pairs outside the finite list
already exceed a larger fixed bound. These statements persist under a
small Hausdorff perturbation, with a strict gap to all pairs outside
$\mathcal E$.

The closest displacement vector between two disjoint convex bodies is
unique: if two different displacement vectors both minimized its norm,
their midpoint would have smaller norm and would still be a displacement.
For that vector, each endpoint is a support point in its normal direction.
Strict convexity makes those support points unique. Rotate the segment
to $[0,g]$ and use a common transverse coordinate. The facing graphs are
$(-\psi_0(u),u)$ and $(g+\psi_1(v),v)$, with
$\psi_b(0)=\psi_b'(0)=0$ and $\psi_b''(0)=\kappa_b>0$.
The distance Hessian at its minimum is
\begin{equation}\label{eq:g-single-hessian}
 \frac1g\begin{pmatrix}1+g\kappa_0&-1\\-1&1+g\kappa_1\end{pmatrix}.
\end{equation}
Its determinant and its diagonal entries are positive. The implicit
function theorem gives smooth contact points and distance. Compactness
keeps both eigenvalues uniformly positive and all chosen charts uniform.

A reference minimum segment cannot meet a third obstacle in its open
interior: the initial part would give a smaller distance from its first
obstacle to that third obstacle. The closed segment also stays a positive
distance from every other obstacle, since its endpoints belong only to
its own two disjoint bodies and only finitely many other bodies are
nearby. Small contact patches and small geometric perturbations preserve
this clearance. Strict convexity makes each connecting chord leave the
first obstacle and enter the second transversely, with no earlier hit.
Thus these local chords are physical first hits.

At a fixed outgoing obstacle, distinct minimizing channels give distinct
normal outgoing collision states. If two shared a contact and outgoing
normal, their other contacts at the same distance would coincide,
contradicting disjointness. The finite collection of states is therefore
separated. Choose their neighborhoods so small that normal reflection
of a flight arriving along $e$ can belong to a short outgoing channel
only in the neighborhood of $\bar e$. This property is open and uniform.

Every physical free flight joins distinct lifted obstacles: a straight
ray leaving a convex body cannot return to it before another hit. Hence
its length is at least $g_*$. If a sum of $j$ such lengths is at most
$jg_*+\varepsilon_0$, every summand is at most $g_*+\varepsilon_0$.
The gap to the other pairs restricts that summand to $\mathcal E$.
Away from the unique contact patches the distance to that pair's
minimum is bounded below; decrease $\varepsilon_0$ to exclude this
complement. Each short flight is then nearly normal. The separated
outgoing states and specular reflection force every following short
flight to be the reverse channel. Induction uses the same neighborhoods
at each impact and introduces no length-dependent constant. Grazing
states and their regular preimages have phase measure zero.
\end{proof}

The lemma also applies to a specified nonminimal closest chord with
positive clearance, for its selected local itinerary. It need not then
exhaust the whole count event. All original six near-circle channels
satisfy the criterion. For example, the criterion also applies to
arbitrarily eccentric ellipses and to noncongruent obstacles on any
fixed lattice, provided the indicated disjointness and positive
curvature conditions hold; there is no small eccentricity assumption.

For one channel the exact one-flight actions are
\begin{equation}\label{eq:g-length}
 \ell_b(u,v)=\sqrt{(g+\psi_b(u)+\psi_{1-b}(v))^2+(v-u)^2},\quad b=0,1.
\end{equation}
Uniformly on a small fixed box,
\begin{align}\label{eq:g-length-jet}
 \ell_b(u,v)&=g+\frac{c_bu^2+c_{1-b}v^2-2uv}{2g}
                  +O((|u|+|v|)^3),\\
 \ell_b(u,v)-g&\ge\psi_b(u)+\psi_{1-b}(v)
                       \ge c_{\rm loc}(u^2+v^2).
 \label{eq:g-coercive}
\end{align}
The full broken-action Hessian has strictly positive diagonal-dominance
margins: $\kappa_b$ in endpoint rows and $2\kappa_b$ in interior rows
at zero. Every row perturbation involves only neighboring coordinates.
The same small box therefore gives strict convexity at every length.

\section{A parameter-dependent boundary-value and relative-flux lemma}
\label{sec:g-action}
Let $W_{j,e}(\xi,u,v)$ be the stationary value of the sum of
\eqref{eq:g-length}, with $y_0=u,y_j=v$. We construct that value on
one common neighborhood. Endpoint Hessians are denoted by
$\mathsf H$; scalar relative remainders are denoted by $\mathcal R$.

\begin{lemma}[Alternating Jacobi reduction]\label{lem:g-jacobi}
Put $c=\sqrt{c_0c_1}$, $\gamma=\arcosh c$,
$\sigma_i=\sqrt{c_{1-(i\bmod2)}}$, and
$D_j=\diag(\sigma_0,\sigma_j)$. The quadratic endpoint Hessian is
\begin{equation}\label{eq:g-effective}
 \mathsf H_{j,e}=\frac{c\sinh\gamma}{g}D_j^{-1}
 \begin{pmatrix}\coth(j\gamma)&-\csch(j\gamma)\\
 -\csch(j\gamma)&\coth(j\gamma)\end{pmatrix}D_j^{-1}.
\end{equation}
Consequently
\begin{equation}\label{eq:g-ratio}
 d^0_{j,e}:=-W_{uv}(\xi,0,0)
  =\frac{c\sinh\gamma}{g\sigma_0\sigma_j\sinh(j\gamma)},\quad
 \sqrt{\det\mathsf H_{j,e}}=\frac{c\sinh\gamma}{g\sigma_0\sigma_j},\quad
 \frac{d^0_{j,e}}{\sqrt{\det\mathsf H_{j,e}}}=\csch(j\gamma).
\end{equation}
The endpoint eigenvalues have positive uniform upper and lower bounds.
The period-two return multiplier is $e^{2\gamma}$.
\end{lemma}
\begin{proof}
The interior equation is
$y_{i+1}-2c_{i\bmod2}y_i+y_{i-1}=0$. Set $y_i=\sigma_i z_i$.
Since $c_{i\bmod2}\sigma_i/\sigma_{i+1}=c$, the recurrence becomes
$z_{i+1}-2cz_i+z_{i-1}=0$ and the whole quadratic action becomes
\[
 \frac c{2g}\sum_{i=0}^{j-1}
       [c(z_i^2+z_{i+1}^2)-2z_iz_{i+1}].
\]
Its endpoint solution is
$z_i=\sinh((j-i)\gamma)z_0/\sinh(j\gamma)
      +\sinh(i\gamma)z_j/\sinh(j\gamma)$.
The first endpoint derivative is $c(cz_0-z_1)/g$.
The identity
$c\sinh(j\gamma)-\sinh((j-1)\gamma)
 =\sinh\gamma\cosh(j\gamma)$, and its reversed version, give
\eqref{eq:g-effective}. Now use $\coth^2-\csch^2=1$.
The middle matrix has eigenvalues $\tanh(j\gamma/2)$ and
$\coth(j\gamma/2)$. Since $\gamma\ge\gamma_*>0$ and the
$\sigma_i$ have uniform positive bounds, this proves the eigenvalue
bounds. The two-periodic scaling conjugates the two-step recurrence
to the square of $\begin{psmallmatrix}2c&-1\\1&0\end{psmallmatrix}$.
The nonzero one-flight twist $1/g$ identifies position-pair and
collision position--momentum coordinates, so its unstable eigenvalue
is the stated physical return multiplier. In particular the formula
at $j=1$ is exactly \eqref{eq:g-single-hessian}.
\end{proof}

\begin{lemma}[Uniform nonlinear bridge and relative determinant]
\label{lem:g-relative}
Suppose $g$ and the two graph functions in \eqref{eq:g-length} range
over a compact smooth parameter family, with $g,\kappa_0,\kappa_1$
bounded above and bounded away from zero. There is a common endpoint
neighborhood on which the stationary segment exists uniquely and is
smooth. For some fixed $\rho<1$, with $a=|u|+|v|$ and
$w_i=\rho^i+\rho^{j-i}$,
\begin{align}\label{eq:g-bridge}
 y_i^{\rm lin}&=\sigma_i\left[
 \frac{\sinh((j-i)\gamma)}{\sinh(j\gamma)}\frac u{\sigma_0}
 +\frac{\sinh(i\gamma)}{\sinh(j\gamma)}\frac v{\sigma_j}\right],\\
 |y_i|&\le Ca w_i,\qquad |y_i-y_i^{\rm lin}|\le Ca^2 w_i,\notag\\
 W_{j,e}-jg&=\tfrac12(u,v)\mathsf H_{j,e}(u,v)^t+O(a^3),\notag\\
 -W_{uv}(\xi,u,v)&=d^0_{j,e}(\xi)b_{j,e}(\xi,u,v),
 \quad b_{j,e}(\xi,0,0)=1,\quad b_{j,e}=1+O(a)>0.
 \label{eq:g-relative}
\end{align}
Every fixed mixed derivative of $y_i$ and $y_i-y_i^{\rm lin}$ is
bounded by $C_k w_i$ times respectively
$a^{\max(1-|\beta|,0)}$ and $a^{\max(2-|\beta|,0)}$, where
$\beta$ counts endpoint derivatives. The action remainder and $b_{j,e}$
have uniform fixed-order bounds. If both graphs are even, the bridge
correction is $O(a^3w_i)$, the action remainder is $O(a^4)$, and
$b_{j,e}-1=O(a^2)$ is even. In an analytic family these bounds hold
on a common complex neighborhood and are analytic bounds.
\end{lemma}
\begin{proof}
Let $H^0_{\rm int}$ be the quadratic interior Hessian. Its inverse,
for $1\le i\le k<j$, is
\begin{equation}\label{eq:v3-green}
 G_{ik}=\frac{g\sigma_i\sigma_k}{c}
 \frac{\sinh(i\gamma)\sinh((j-k)\gamma)}
      {\sinh\gamma\sinh(j\gamma)},\qquad G_{ki}=G_{ik}.
\end{equation}
Substitution verifies the homogeneous recurrence away from $i=k$,
zero endpoint values and the unit jump at $i=k$. In particular
$|G_{ik}|\le Ce^{-\gamma_*|i-k|}$.
Choose $e^{-\gamma_*}<\rho_0<\rho<1$. A geometric-series sum yields
\begin{equation}\label{eq:v3-weighted-green}
 \sum_k\rho_0^{|i-k|}(\rho^k+\rho^{j-k})\le C_\rho w_i.
\end{equation}
For its $\rho^k$ part, split at $k=i$ and sum the ratios
$\rho_0/\rho$ and $\rho_0\rho$; the other part follows by reversal.
Thus $G$ is bounded independently of $j$ on the norm
$\|z\|_w=\max_i|z_i|/w_i$.

Parameter differentiation of \eqref{eq:v3-green} has the same property.
To see this without a hidden growing dimension, write the hyperbolic
quotient as a bounded factor times $e^{-\gamma(k-i)}$, using factors
$1-e^{-2r\gamma}$ for its distances to the endpoints. Derivatives
of those factors produce $r^m e^{-2r\gamma}$, uniformly bounded for
each fixed $m$; derivatives of the first exponential produce
$|i-k|^m e^{-\gamma|i-k|}\le C_m\rho_0^{|i-k|}$.
All derivatives of the positive $\sigma_i$ are bounded. The same
argument for the endpoint solution gives
$|\partial_\xi^\alpha y_i^{\rm lin}|\le C_\alpha a w_i$.
Consequently every fixed parameter derivative of $G$ is bounded in
the weighted norm, with no $j$ factor.

The interior gradient is $H^0_{\rm int}(y-y^{\rm lin})+
\mathcal N_\xi(y)$, including the fixed endpoint coordinates in the
local remainder. Its $i$th component depends on only three neighboring
coordinates and vanishes quadratically there. Adjacent weights have
bounded ratios. On $\|y\|_w\le C_0a$, its weighted norm is
$O(a^2)$ and its derivative has norm $O(a)$. The map
$y\mapsto y^{\rm lin}-G\mathcal N_\xi(y)$ preserves this ball
and has Lipschitz constant at most $1/2$ when $a$ is uniformly small.
This proves existence and the first estimates. The diagonal-dominance
bound after \eqref{eq:g-coercive} gives uniqueness throughout the
contact box, not just inside the weighted ball.

For clarity, the derivative induction takes place in these same
weighted spaces. One derivative of the fixed-point equation gives
\[
 (I+G D_y\mathcal N_\xi)\,D y
     =D y^{\rm lin}-(D G)\mathcal N_\xi
                         -G D_{\rm explicit}\mathcal N_\xi.
\]
Its left inverse has norm at most two. At order $k$, the only term
containing an order-$k$ derivative of $y$ is the same left operator;
every other term contains derivatives of $G$, explicit parameter
or endpoint derivatives of the local remainder, and lower derivatives
of $y$. The local chain rule and bounded neighboring-weight ratios
bound their products in $\|\cdot\|_w$. Induction proves the estimates
at order $k$. Parameter derivatives preserve vanishing at zero;
Taylor's integral formula in $(u,v)$ therefore gives precisely the
powers of $a$ stated in the lemma, including the correction of order two.

For every $p>0$, $\sum_i w_i^p\le C_p$. Sum the cubic action
remainders using this inequality. The linear bridge is stationary for
the quadratic action with endpoints fixed, so its change on substituting
the nonlinear bridge is $O(a^4)$. This proves the action estimate.
The same summability and the derivative induction prove all its
fixed-order bounds.

For the relative estimate put
$b_i=-\ell_{i\bmod2,uv}(y_i,y_{i+1})$ and
$H_{\rm int}=H^0_{\rm int}+E$. Schur complementation gives
\begin{equation}\label{eq:v3-cofactor}
              -W_{uv}=\frac{\prod_{i=0}^{j-1}b_i}{\det H_{\rm int}}.
\end{equation}
Indeed, the endpoint-to-interior entries are $-b_0,-b_{j-1}$ and
the corner entry of the inverse tridiagonal interior matrix is
$b_1\cdots b_{j-2}/\det H_{\rm int}$. At $j=1$ the determinant
is empty and equals one. The local expansions and endpoint weights give
\[
 \sum_i|\log(gb_i)|\le Ca,\qquad
 \|E\|_1\le\sum_{i,k}|E_{ik}|\le Ca,\qquad
 \|(H^0_{\rm int})^{-1}\|\le C.
\]
The trace-norm bound uses tridiagonality; each matrix unit has trace
norm one. Take the endpoint neighborhood so small that
$\|(H^0_{\rm int})^{-1}E\|<1/2$. Then
\begin{equation}\label{eq:v3-logdet}
 \log\frac{\det H_{\rm int}}{\det H^0_{\rm int}}
 =\sum_{m\ge1}\frac{(-1)^{m+1}}m
       \tr[((H^0_{\rm int})^{-1}E)^m]=O(a).
\end{equation}
The absolute trace is bounded by one trace norm times $m-1$ operator
norms. Summing therefore costs $C\|E\|_1$, not the matrix dimension.
Divide \eqref{eq:v3-cofactor} by its value at zero and subtract
logarithms to obtain \eqref{eq:g-relative}.

Every differentiated $E$ has uniformly bounded trace norm by the
already proved coordinate bounds. The quadratic Hessian and each of
its parameter derivatives are uniformly banded with bounded entries;
hence derivatives of its inverse have bounded operator norm, by
$\partial H^{-1}=-H^{-1}(\partial H)H^{-1}$ and induction.
Termwise differentiation of \eqref{eq:v3-logdet} is justified by
geometric convergence even after its fixed polynomial factors in $m$.
Every term has at least one trace-class factor. The differentiated sum
of $\log(gb_i)$ is controlled by the same endpoint weights. This proves
all derivative bounds for the logarithmic ratio and its exponential.
Positive twists and a positive interior determinant give positivity.

If the graphs are even, the action expansion is even, its nonlinear
gradient is cubic, and its Hessian and twist errors are quadratic.
The preceding estimates then give the improved orders. Uniqueness
implies simultaneous reversal of every coordinate, proving evenness.
For a jointly analytic family all local functions extend to a common
small complex neighborhood of the fixed compact real parameter set.
Shrink it to retain the strict exponential and inverse-norm margins.
The same contraction is uniform there. The determinant series remains
convergent, and Cauchy estimates on a smaller neighborhood prove the
analytic assertion. No assertion about analytic extension is used for
a merely smooth family.
\end{proof}
